% Options for packages loaded elsewhere
\PassOptionsToPackage{unicode}{hyperref}
\PassOptionsToPackage{hyphens}{url}
\documentclass[
  a4paper,
]{article}
\usepackage{xcolor}
\usepackage[margin=25mm]{geometry}
\usepackage{amsmath,amssymb}
\setcounter{secnumdepth}{-\maxdimen} % remove section numbering
\usepackage{iftex}
\ifPDFTeX
  \usepackage[T1]{fontenc}
  \usepackage[utf8]{inputenc}
  \usepackage{textcomp} % provide euro and other symbols
\else % if luatex or xetex
  \usepackage{unicode-math} % this also loads fontspec
  \defaultfontfeatures{Scale=MatchLowercase}
  \defaultfontfeatures[\rmfamily]{Ligatures=TeX,Scale=1}
\fi
\usepackage{lmodern}
\ifPDFTeX\else
  % xetex/luatex font selection
\fi
% Use upquote if available, for straight quotes in verbatim environments
\IfFileExists{upquote.sty}{\usepackage{upquote}}{}
\IfFileExists{microtype.sty}{% use microtype if available
  \usepackage[]{microtype}
  \UseMicrotypeSet[protrusion]{basicmath} % disable protrusion for tt fonts
}{}
\makeatletter
\@ifundefined{KOMAClassName}{% if non-KOMA class
  \IfFileExists{parskip.sty}{%
    \usepackage{parskip}
  }{% else
    \setlength{\parindent}{0pt}
    \setlength{\parskip}{6pt plus 2pt minus 1pt}}
}{% if KOMA class
  \KOMAoptions{parskip=half}}
\makeatother
\setlength{\emergencystretch}{3em} % prevent overfull lines
\providecommand{\tightlist}{%
  \setlength{\itemsep}{0pt}\setlength{\parskip}{0pt}}
\usepackage{bookmark}
\IfFileExists{xurl.sty}{\usepackage{xurl}}{} % add URL line breaks if available
\urlstyle{same}
\hypersetup{
  pdftitle={Designing a Research Project to Explore the Emergence of Spacetime{,} Particles{,} and Interactions from States and Relations --- A Framework of Preliminary Design and Feasibility Verification for Foundational-Physics Model Exploration under Uncertainty},
  pdfauthor={Noriaki Kihara (WF System Co., Ltd.) ORCID 0009-0004-6753-4020},
  hidelinks,
  pdfcreator={LaTeX via pandoc}}

\title{Designing a Research Project to Explore the Emergence of
Spacetime, Particles, and Interactions from States and Relations --- A
Framework of Preliminary Design and Feasibility Verification for
Foundational-Physics Model Exploration under Uncertainty}
\author{Noriaki Kihara (WF System Co., Ltd.) ORCID 0009-0004-6753-4020}
\date{September 20, 2026 Version DOI 10.5281/zenodo.22851945}

\begin{document}
\maketitle

\textbf{Author:} Noriaki Kihara\\
\textbf{ORCID:} 0009-0004-6753-4020\\
\textbf{Version DOI:} 10.5281/zenodo.22851945\\
\textbf{Concept DOI:} 10.5281/zenodo.22851944\\
\textbf{Version:} v1.0\\
\textbf{Date:} 2026-09-20

\begin{center}\rule{0.5\linewidth}{0.5pt}\end{center}

\subsection{Abstract}\label{abstract}

The Standard Model and General Relativity are extremely successful
within their respective domains of application. At the same time, the
problem of describing quantum theory and gravity naturally from a single
underlying structure remains unsolved: despite decades of work on
several major research programs, including superstring theory, loop
quantum gravity, asymptotic safety, and holography, no completed picture
uniquely recovers real low-energy physics, a quantum description of
spacetime, the cosmological background, and observables at the same time
{[}4--9{]}. Moreover, while the Standard Model possesses high gauge
symmetry and conservation structure, the description of concrete
physical states requires multiple experimentally determined parameters
{[}10{]}.

The author has been exploring whether spacetime, fields, particles,
mass, charge, distance, and time --- rather than being given from the
start --- can be generated spontaneously from a smaller number of
``states'' and ``relations'' alone. However, when individual candidate
equations and numerical experiments are put first, the research
direction tends to shift whenever a locally interesting phenomenon
appears, and the exploration proceeds while the project's goal image ---
``what structure must be generated for the research objective to count
as achieved'' --- remains vague.

This paper presents no completed theory. Instead, applying
project-management research on exploratory projects under high
uncertainty {[}1--3{]}, together with the author's own working
heuristics from IT consulting and systems design, it redesigns this
research as a long-term exploratory project. Concretely, it defines an
iterative research structure of philosophical principles, a provisional
goal image, preliminary design, feasibility verification, numerical
experiments, and design revision.

The goal of this research is not to invent a new physical law and
replace existing theories. It is to explore whether there exists an
underlying structure or interpretation that can explain the Standard
Model and General Relativity --- with their successes preserved --- from
fewer assumptions and fewer arbitrary parameters.

\textbf{Keywords:} foundational physics, quantum gravity, emergent
spacetime, states and relations, exploratory project, project
management, preliminary design, feasibility verification

\begin{center}\rule{0.5\linewidth}{0.5pt}\end{center}

\subsection{1. Motivation}\label{motivation}

The starting point of this research is the following two questions.

First: given that the Standard Model and General Relativity are each
successful to very high precision, why is it difficult to derive both
naturally from the same underlying structure?

General Relativity treats spacetime geometry itself as a dynamical
degree of freedom. The Standard Model, on the other hand, describes the
strong, weak, and electromagnetic interactions on the basis of quantum
field theory and gauge symmetry. At low energies, General Relativity can
be treated quantum mechanically as an effective field theory, but no
completed form of a fundamental quantum theory of gravity valid up to
the Planck regime has been obtained {[}4,5{]}.

Second: while the Standard Model possesses high symmetry and many
conservation structures, its concrete description requires multiple
experimental parameters. This does not mean the Standard Model has
failed. Rather, this research considers the possibility that the
existing theories are highly accurate effective descriptions, behind
which there exists a foundational layer that generates the same
structures from a smaller number of principles.

String theory has important achievements as a quantum framework that
includes gravity, and gauge--gravity duality and holography give strong
hints about the relation between spacetime and information. However,
unresolved problems remain concerning background independence,
non-perturbative formulation, the selection of realistic vacua, and
application to cosmological spacetimes {[}6,7{]}. Loop quantum gravity,
asymptotic safety, discrete spacetime models, and others have likewise
shown important progress from different angles, but no consensus has
been reached on a single completed theory {[}4,8,9{]}.

This research does not deny these existing approaches. Rather, it treats
the successes of the existing theories as strong constraints that must
ultimately be reproduced.

\begin{center}\rule{0.5\linewidth}{0.5pt}\end{center}

\subsection{2. Position and Scope as an Independent
Researcher}\label{position-and-scope-as-an-independent-researcher}

The author is an independent researcher, not affiliated with a
university, research institute, or academic society specializing in
quantum gravity, particle theory, or general relativity.

This position has two sides.

On one hand, the author is not directly bound by the research programs,
particular theoretical lineages, or community-internal problem settings
usually required inside established research fields, and can explore by
combining concepts from different fields and many hypotheses relatively
freely. This makes it possible to test the strong hypothesis of removing
the existing concepts of spacetime, fields, particles, mass, and charge
themselves from the initial premises and reconstructing them from a
smaller number of states and relations.

On the other hand, research resources are limited. Carrying out
large-scale collaboration, dedicated computational resources, long-term
rigorous mathematical analysis, and complete connection proofs to
existing theories simultaneously is not realistic. For this reason, the
outputs of this research primarily target the level of theoretical
proposals, thought experiments, numerical experiments, and reproducible
preliminary designs, and this project alone does not aim to derive a
complete rigorous connection to the Standard Model or General
Relativity.

This limitation does not mean neglecting consistency with existing
theories. Quite the opposite. In this research, the Standard Model and
General Relativity are treated not as ``objects to be replaced'' but as
``known successes to be reproduced.''

The purpose of this research is therefore not

\[
\boxed{\text{to invent new physics}}
\]

but

\[
\boxed{
\text{to explore structures that can explain the known Standard Model and relativity}\
\text{from fewer assumptions and fewer arbitrary parameters.}
}
\]

The structure presented in this paper is meant to fix the direction of
the research project toward this purpose; it is not the claim of a
completed physical theory.

\begin{center}\rule{0.5\linewidth}{0.5pt}\end{center}

\subsection{3. Why the Research Project Itself Is a Design
Object}\label{why-the-research-project-itself-is-a-design-object}

In this research, the author has so far repeated an exploration cycle of
building individual candidate equations, running numerical experiments,
and analyzing the phenomena obtained. This method is heuristically
effective, but it has one problem.

When candidate equations yield interesting phenomena --- singularities,
closure structures, metastable states, exponential growth, periodic
structures --- those phenomena easily start to look as if they
themselves were the research objective.

However,

\[
\boxed{
\text{an interesting numerical phenomenon}
\neq
\text{the foundational structure actually being sought.}
}
\]

Project-management research has likewise pointed out that projects whose
exploration methods and final goals are vague require learning,
candidate selection, and iterative design revision, not merely fixed
plan execution {[}1--3{]}.

Turner and Cochrane classified projects by separating the clarity of
goals from the clarity of the methods of achieving them {[}1{]}. Pich,
Loch, and De Meyer argued for the need to use instructionism, learning,
and selectionism according to the kinds of uncertainty, ambiguity, and
complexity involved {[}2{]}. Lenfle showed that exploratory projects, in
which neither the technology nor the requirements are fixed at the
start, need management methods different from those of ordinary
development projects {[}3{]}.

In this sense, this research is close to a typical exploratory project.

The following iterative structure is therefore adopted:

\[
\boxed{
\text{philosophical principles}
\rightarrow
\text{goal image}
\rightarrow
\text{preliminary design}
\rightarrow
\text{feasibility verification}
\rightarrow
\text{numerical experiments}
\rightarrow
\text{revision of design and goal image}
}
\]

What matters is that the goal image itself is not a fixed specification.
If numerical experiments or mathematical examination show it to be
impossible or unnatural, the goal image itself is revised.

\begin{center}\rule{0.5\linewidth}{0.5pt}\end{center}

\subsection{4. First-Principle Candidate: Nameless Distinguishable
States}\label{first-principle-candidate-nameless-distinguishable-states}

This research keeps what exists initially to a minimum.

The first-principle candidate is expressed as follows:

\[
\boxed{
\text{There exist states that have no names, yet are mutually distinguishable and countable.}
}
\]

For convenience, the state set is written

\[
\mathcal Z_n=\{z_{1,n},z_{2,n},\ldots,z_{N_n,n}\}.
\]

Here the subscripts are computational bookkeeping identifiers, not
physical ``names.'' States are not endowed in advance with attributes
such as electron, photon, mass, charge, spatial position, or time
coordinate.

\begin{center}\rule{0.5\linewidth}{0.5pt}\end{center}

\subsection{5. Non-direct Readability and Interaction =
Readout}\label{non-direct-readability-and-interaction-readout}

This research does not assume that a state itself can be read directly.

Likewise, it does not assume that a relation between states can itself
be read directly.

Physical readout is taken to become possible only as a product with
other states or relations.

Conceptually,

\[
\boxed{
\text{a state alone: unreadable}
}
\]

\[
\boxed{
\text{a relation alone: unreadable}
}
\]

\[
\boxed{
\text{product-type operations of states and relations: readable.}
}
\]

In this position, observation and interaction cannot be separated in
principle.

To read anything out, an operation with other states and relations is
required, and that operation itself produces the next state.

This research therefore takes

\[
\boxed{
\text{observation}\equiv\text{readout accompanied by interaction.}
}
\]

\begin{center}\rule{0.5\linewidth}{0.5pt}\end{center}

\subsection{6. Self-interaction Is Not
Excluded}\label{self-interaction-is-not-excluded}

Since states have no physical names, there is initially no ground for
excluding only oneself from the interaction partners.

Hence, writing the unknown two-state map as

\[
f(z_i,z_j),
\]

the minimal skeleton producing each state's next state is taken as the
candidate

\[
\boxed{
 z_{i,n+1}
 =
 \sum_{j=1}^{N_n} f(z_{i,n},z_{j,n}).
}
\]

The sum includes

\[
j=i
\]

as well.

However, this paper does not decide the concrete

\[
f(z_i,z_j).
\]

This is precisely the main object to be explored in the coming
preliminary design and feasibility verification.

\begin{center}\rule{0.5\linewidth}{0.5pt}\end{center}

\subsection{7. The Project's First Goal
Image}\label{the-projects-first-goal-image}

The following is neither an axiom nor a completed requirements
specification.

It is the first goal image --- what the author currently hopes will
ultimately appear if structures resembling known physics can be
generated from a small number of states and relations alone.

\subsubsection{7.1 Global conserved
quantity}\label{global-conserved-quantity}

The system as a whole is expected to have a conserved quantity

\[
C_{\mathrm{tot}}
\]

that may later be read out as corresponding to mass or energy.

For example, a structure like

\[
\boxed{
\sum_{i=1}^{N_n} C(z_{i,n})=C_{\mathrm{tot}}.
}
\]

The concrete form of the conserved quantity is undetermined.

\subsubsection{7.2 Growth and decrease of the number of
states}\label{growth-and-decrease-of-the-number-of-states}

The number of states

\[
N_n
\]

is not assumed fixed.

Both

\[
N_{n+1}>N_n
\]

and

\[
N_{n+1}<N_n
\]

are allowed.

A structure in which the global conserved quantity is maintained even as
the number of states changes is desirable.

\subsubsection{7.3 Extension arising from
relations}\label{extension-arising-from-relations}

Between states, multiple relational quantities

\[
R_{ij}^{(\mu)}
\]

are expected to form, some of which may later be read out as distance or
spatial extension.

The direction \(\mu\) is not given as a spatial axis from the start.

Multiple directions are expected to form spontaneously from
interactions.

\subsubsection{7.4 Initial highly symmetric
state}\label{initial-highly-symmetric-state}

Initially, a highly symmetric state is envisaged in which the
all-direction total of the relational quantities that look like distance
or extension is arbitrarily close to zero:

\[
\boxed{
L_\mu(0)\simeq 0.
}
\]

However, exact zero is not fixed from outside.

\subsubsection{7.5 Spontaneous symmetry breaking and exponential
extension}\label{spontaneous-symmetry-breaking-and-exponential-extension}

If the highly symmetric state is not a stable fixed point but an
unstable one, a small asymmetry

\[
\epsilon
\]

may be amplified as

\[
\epsilon\rightarrow e^{\lambda n}\epsilon,
\]

producing extension without any externally specified direction.

For directions that look spatial, the appearance, over some period, of
exponential growth like

\[
\boxed{
L_s(n)\sim e^{Hn}
}
\]

is taken as one element of the goal image.

\subsubsection{7.6 Temporal direction}\label{temporal-direction}

By contrast, one relational direction is expected to differ from the
spatial ones: not condensing locally, but increasing monotonically as
the accumulation of causal processing.

For example, an approximately linear increase like

\[
\boxed{
L_t(n)\sim n
}
\]

may later be read out as time.

Rather than giving time as a continuous coordinate from the start, it is
desirable to generate it as the order, or accumulated amount, of
discrete causal updates.

\subsubsection{7.7 Upper bound of extension and phase
transition}\label{upper-bound-of-extension-and-phase-transition}

Spatial exponential growth is not expected to continue without limit; it
is expected to reach a ceiling originating from the number of states,
the number of relations, the conserved quantity, or the total of
relational quantities, and to undergo another phase transition.

As a result,

\[
\boxed{
\text{growth of wide-area relations}
\rightarrow
\text{condensation of local relations}
}
\]

is considered as a possible outcome.

\subsubsection{7.8 Condensates and
particles}\label{condensates-and-particles}

If locally stable relational structures form, they are read out
afterwards as particles or condensates.

Particles are not a basic input.

Condensation is expected to occur in the spatial directions, while the
same kind of condensation is expected not to occur in the temporal
direction.

\subsubsection{7.9 Reactions between
condensates}\label{reactions-between-condensates}

In interactions between condensates, the conserved quantity is expected
to appear not as unlimited fresh creation, but as

\[
\boxed{
\text{transmission}+\text{partition}+\text{redistribution}.
}
\]

Furthermore, it is a long-term goal that both gravitational and
electromagnetic interactions later appear as different readouts of the
same basic states and the same relational operations.

\begin{center}\rule{0.5\linewidth}{0.5pt}\end{center}

\subsection{8. Performance Targets: Cutoff Criteria for Evaluating
Candidate
Models}\label{performance-targets-cutoff-criteria-for-evaluating-candidate-models}

The following are also not axioms at this stage; they are performance
targets for evaluating preliminary designs.

\subsubsection{8.1 Boundedness}\label{boundedness}

The most important performance target is boundedness.

From finite input,

\begin{itemize}
\tightlist
\item
  state values
\item
  relational quantities
\item
  the number of states
\item
  the number of relations
\item
  the operational steps needed for one interaction
\end{itemize}

are expected not to diverge without limit.

Conceptually,

\[
\boxed{
\text{finite input}
\rightarrow
\text{finite operations}
\rightarrow
\text{finite output.}
}
\]

This differs from assuming that the universe as a whole has a
permanently fixed maximum number of states. The requirement is that at
any finite processing stage, the computational objects and physical
quantities are finite.

\subsubsection{8.2 Discrete causality}\label{discrete-causality}

An interaction has the definite order

\[
\text{input}
\rightarrow
\text{operation}
\rightarrow
\text{output}.
\]

The basic description is therefore expected to be

\[
\boxed{
\mathcal Z_n\rightarrow\mathcal Z_{n+1}.
}
\]

In this position, the elementary process is essentially discrete, and
there is no need to postulate intermediate states

\[
n+\epsilon
\]

as physical reality.

Differentiable continuous equations are considered to appear, if at all,
only as effective descriptions in a coarse-grained limit of many
discrete updates.

\subsubsection{8.3 Resolution floor}\label{resolution-floor}

This research does not assume that the infinitesimal and zero are always
physically distinguishable.

Readout is performed not by direct measurement of a single state, but by
product-sum operations with other states and relations.

Conceptually, a finite-range product-sum readout like

\[
Y_{AB}
=
\sum_k P(R_{Ak},R_{Bk})
\]

is considered.

If this operation acts as a finite-bandwidth integrating operation, the
readout itself may possess an effective lower bound of resolution.

If the difference of two states is below the readout resolution, it is
treated as

\[
\boxed{
\text{an unresolvable nonzero}\equiv\text{physically zero.}
}
\]

What matters is that this resolution is not given from the start as a
fixed constant. It is expected that the two interacting relational
structures themselves determine the resolution, as in

\[
\boxed{
\epsilon_{AB}=\mathcal E(R_A,R_B).
}
\]

Note that a low-pass characteristic does not follow automatically from
``being a product-sum'' alone. Whether high-frequency components are
actually suppressed to yield a finite bandwidth is an item to be
verified for each concrete candidate map.

\begin{center}\rule{0.5\linewidth}{0.5pt}\end{center}

\subsection{9. The Position of Quantum
Behavior}\label{the-position-of-quantum-behavior}

This research does not place quantum behavior as a first axiom.

If, from

\[
\boxed{
\text{discrete causal processing}
+
\text{finite resolution}
+
\text{product-sum relational readout,}
}
\]

distinguishable observational outcomes become naturally discretized,
then quantum behavior may be explainable as an emergent property.

Accordingly, quantum behavior is treated not as

\[
\boxed{
\text{quantum behavior}
=
\text{an axiom}
}
\]

but as

\[
\boxed{
\text{quantum behavior}
=
\text{a phenomenon to be derived from discrete causality and finite resolution.}
}
\]

This is not a conclusion proven in this paper; it is a subject of future
verification.

\begin{center}\rule{0.5\linewidth}{0.5pt}\end{center}

\subsection{10. Connection Policy toward the Standard Model and General
Relativity}\label{connection-policy-toward-the-standard-model-and-general-relativity}

This research does not start from denying or replacing the Standard
Model and General Relativity.

Rather, if a candidate model matures sufficiently, it must recover at
least the following two kinds of limits.

First, for gravity: reproducing the Newtonian distance dependence in the
weak-field, low-velocity limit, and more generally yielding readouts
consistent with the geometric description of General Relativity.

Second, for non-gravitational interactions: conservation laws,
symmetries, and interaction structures corresponding to the gauge
interactions of the Standard Model being read out as effective
descriptions.

At the current stage of the project, however, rigorous derivation of
these is not an immediate goal.

First,

\[
\boxed{
\text{does there exist a finite dynamics that can generate these structures from states and relations alone?}
}
\]

is examined at the level of feasibility.

In existing quantum-gravity research as well, the emergent viewpoint ---
classical spacetime and gravitons appearing in suitable limits from more
basic short-distance structure --- is an important research theme
{[}8{]}. The holographic principle suggests a deep relation between
geometry and information {[}7{]}. This research does not adopt these
specific theories, but the exploratory direction itself --- not taking
known spacetime and particles as basic entities from the start --- is
not isolated from the problem-consciousness of contemporary
quantum-gravity research.

\begin{center}\rule{0.5\linewidth}{0.5pt}\end{center}

\subsection{11. Preliminary Design and Feasibility
Verification}\label{preliminary-design-and-feasibility-verification}

At the current stage, preliminary design and feasibility verification
are needed before proceeding to basic design.

The reason is simple.

There is no guarantee that a map satisfying all the goal images listed
in this paper exists.

Moreover, even if one exists, it may not be discoverable within
realistic research resources.

Therefore, instead of fixing one map as the correct answer, multiple
candidates

\[
f_1,
\ f_2,
\ f_3,
\ \ldots
\]

are constructed and evaluated in stages.

The first feasibility verification is limited to a small number of
conditions, for example:

\begin{enumerate}
\def\labelenumi{\arabic{enumi}.}
\tightlist
\item
  Can a global conserved quantity be maintained?
\item
  Can all-state interaction including self-interaction be defined?
\item
  Can the conservation structure be maintained as the number of states
  grows or shrinks?
\item
  Do states, relations, and operations remain finite?
\item
  Are relational directions generated without externally supplied
  spatial axes or particle species?
\end{enumerate}

Candidates that fail these are not extended to the later stages of
cosmic expansion, condensation, gravity, and electromagnetism.

When a candidate fails, instead of immediately adding auxiliary
assumptions to keep it alive, the question

\[
\boxed{
\text{is the candidate map wrong?}
}
\]

or

\[
\boxed{
\text{is the goal image itself wrong?}
}
\]

is examined separately.

This iteration is the central method of this research.

\begin{center}\rule{0.5\linewidth}{0.5pt}\end{center}

\subsection{12. Research Prohibitions}\label{research-prohibitions}

While preserving the freedom of exploration, this project avoids, as far
as possible, hidden inputs such as the following.

\begin{itemize}
\tightlist
\item
  Giving particle-species names to initial states.
\item
  Separating basic waves for gravity and for electromagnetism from the
  start.
\item
  Fixing spatial or temporal axes externally.
\item
  Pre-specifying the network of interaction partners without reason.
\item
  Excluding only self-interaction without reason.
\item
  Adding normalization or clipping solely to obtain desirable results.
\item
  Repairing a diverging candidate with a saturation function that has no
  physical grounding.
\item
  Interpreting an observed numerical phenomenon as a physical discovery
  without verification of correspondence with known physics.
\end{itemize}

These are not absolute prohibitions. If their necessity is derived, they
may be introduced. In that case, however, they must be stated explicitly
as new assumptions.

\begin{center}\rule{0.5\linewidth}{0.5pt}\end{center}

\subsection{13. Success and Failure
Conditions}\label{success-and-failure-conditions}

Success for this project does not mean only obtaining a completed theory
of everything.

For example, any of the following counts as an outcome:

\begin{itemize}
\tightlist
\item
  Showing that a minimal map exists that simultaneously satisfies part
  of the goal image.
\item
  Showing that certain requirements are mathematically incompatible with
  each other.
\item
  Showing that a specific known physical structure can emerge from
  states and relations alone.
\item
  Conversely, showing that additional basic degrees of freedom are
  indispensable.
\item
  Establishing effective exploration procedures and rejection criteria.
\end{itemize}

Therefore,

\[
\boxed{
\text{the discovery of infeasibility is also a research outcome.}
}
\]

\begin{center}\rule{0.5\linewidth}{0.5pt}\end{center}

\subsection{14. Significance of This
Paper}\label{significance-of-this-paper}

This paper contains no proof of a new physical law.

Its purpose is to clearly separate what has tended to be mixed together:

\[
\boxed{
\text{philosophical principles}
}
\]

\[
\boxed{
\text{axiom candidates}
}
\]

\[
\boxed{
\text{the project's goal image}
}
\]

\[
\boxed{
\text{assumptions of preliminary design}
}
\]

\[
\boxed{
\text{phenomena obtained in numerical experiments}
}
\]

\[
\boxed{
\text{conditions for comparison with known physics.}
}
\]

In foundational-physics exploration, obtaining a locally interesting
phenomenon and that structure explaining real physics are not the same
thing.

This project therefore emphasizes

\[
\boxed{
\text{preserving not only what was discovered, but why that computation was performed.}
}
\]

This is also a research-management principle for maintaining
reproducibility and direction in independent research that must sustain
long-term exploration with limited resources.

\begin{center}\rule{0.5\linewidth}{0.5pt}\end{center}

\subsection{15. Conclusion}\label{conclusion}

The immediate purpose of this research is not to present a completed
equation of the universe.

More restrictedly, it is to examine

\[
\boxed{
\text{whether there exists a finite computational system that can spontaneously generate,}
\text{from nameless distinguishable states and relations alone,}\
\text{spacetime, particle, and interaction structures resembling known physics.}
}
\]

This paper adopted no specific two-body interaction map.

Instead, it defined the research process itself:

\[
\boxed{
\text{philosophical principles}
\rightarrow
\text{provisional goal image}
\rightarrow
\text{preliminary design}
\rightarrow
\text{feasibility verification}
\rightarrow
\text{numerical experiments}
\rightarrow
\text{revision.}
}
\]

As an independent researcher, the author does not deny existing
theories, but explores whether structures exist that can explain their
successes from fewer assumptions.

In that process, it is accepted that candidate maps will be rejected,
that the goal image itself will be revised, and that the original
conception may ultimately turn out to be infeasible.

This paper is the planning document of that exploration --- written so
that the search neither drifts nor is fixed prematurely into a single
theory.

\begin{center}\rule{0.5\linewidth}{0.5pt}\end{center}

\subsection{References}\label{references}

{[}1{]} J. R. Turner and R. A. Cochrane, ``Goals-and-methods matrix:
coping with projects with ill defined goals and/or methods of achieving
them,'' \emph{International Journal of Project Management}, 11(2),
93--102 (1993). DOI: 10.1016/0263-7863(93)90017-H.

{[}2{]} M. T. Pich, C. H. Loch, and A. De Meyer, ``On Uncertainty,
Ambiguity, and Complexity in Project Management,'' \emph{Management
Science}, 48(8), 1008--1023 (2002). DOI: 10.1287/mnsc.48.8.1008.163.

{[}3{]} S. Lenfle, ``Exploration and project management,''
\emph{International Journal of Project Management}, 26(5), 469--478
(2008). DOI: 10.1016/j.ijproman.2008.05.017.

{[}4{]} S. Carlip, D.-W. Chiou, W.-T. Ni, and R. Woodard, ``Quantum
Gravity: A Brief History of Ideas and Some Prospects,''
\emph{International Journal of Modern Physics D}, 24 (2015). DOI:
10.1142/S0218271815300281. arXiv:1507.08194.

{[}5{]} J. F. Donoghue, ``Quantum General Relativity and Effective Field
Theory,'' arXiv:2211.09902 (2022).

{[}6{]} M. Blau and S. Theisen, ``String theory as a theory of quantum
gravity: a status report,'' \emph{General Relativity and Gravitation},
41, 743--755 (2009). DOI: 10.1007/s10714-008-0752-z.

{[}7{]} R. Bousso, ``The Holographic Principle,'' \emph{Reviews of
Modern Physics}, 74, 825--874 (2002). DOI: 10.1103/RevModPhys.74.825.

{[}8{]} A. Ashtekar, M. Reuter, and C. Rovelli, ``From General
Relativity to Quantum Gravity,'' arXiv:1408.4336 (2014); subsequently
published as a chapter in \emph{General Relativity and Gravitation: A
Centennial Perspective}.

{[}9{]} L. Smolin, ``How far are we from the quantum theory of
gravity?'' arXiv:hep-th/0303185 (2003).

{[}10{]} Particle Data Group, \emph{Review of Particle Physics}, current
Standard Model and related reviews; see F. Takahashi et al.~(Particle
Data Group), \emph{International Journal of Modern Physics A} 41,
2630011 (2026).

\begin{center}\rule{0.5\linewidth}{0.5pt}\end{center}

\subsection{Note: Citation Policy of This
Paper}\label{note-citation-policy-of-this-paper}

Because the purpose of this paper is to fix the direction of the
research project, the author's own previously published papers are not
used as grounds. There is no self-citation, and external citations are
limited to (i) management methods for exploratory projects under high
uncertainty, and (ii) reviews organizing the current position of the
Standard Model, General Relativity, and quantum-gravity research.

\end{document}
